\documentclass[11pt]{article}
\usepackage[margin=1in]{geometry}
\usepackage{amsmath,amssymb}
\usepackage[T1]{fontenc}
\usepackage[hidelinks]{hyperref}

\newcommand{\Hh}{\mathcal H}
\newcommand{\E}{\mathbb E}

\title{Literature-Implied Status Note for arXiv:2004.00891\\
Trace-Norm Consistency Without Reflexivity}
\author{Run fa\_banach\_001, agent\_lane\_08}
\date{11 August 2026}

\begin{document}
\maketitle

\section*{Status}

\textbf{literature\_implied\_answer (full trace-class conclusion).}
This is an identification of a known 1972 theorem, not a new theorem.

\section*{Source question}

Mollenhauer--Klus--Sch\"utte--Koltai,
\emph{Kernel Autocovariance Operators of Stationary Processes: Estimation
and Convergence}, arXiv:2004.00891v2, Remark~11 on PDF page~9, observes that
the trace class $S_1(\Hh)$ is nonreflexive when $\Hh$ is infinite-dimensional.
It asks whether reflexivity is necessary for the almost-sure norm convergence
in the Banach-valued Birkhoff theorem, with convergence of the empirical
kernel autocovariance in trace norm as the motivating unresolved case.

\section*{Supporting literature}

Pop-Stojanovic, \emph{On ergodic theorem for a Banach valued random
sequence}, J. Austral. Math. Soc. \textbf{13} (1972), 501--507,
Corollary~2 on journal page~505 (proof on page~506; PDF pages~5--6), states
strong almost-sure convergence of the averages of a Bochner-integrable
strictly stationary Banach-valued sequence to its invariant conditional
expectation.  No reflexivity assumption is imposed.  In an ergodic system the
invariant conditional expectation is constant and equals the Bochner
expectation.

\section*{Identification}

Fix an integer lag $\eta$ and set
\[
 F_t=\varphi(X_{t+\eta})\otimes\varphi(X_t)\in S_1(\Hh).
\]
The rank-one trace-norm identity and Cauchy--Schwarz give
\[
 \E\|F_0\|_{S_1}
 =\E\bigl[\|\varphi(X_\eta)\|\,\|\varphi(X_0)\|\bigr]
 \leq
 \bigl(\E\|\varphi(X_\eta)\|^2\bigr)^{1/2}
 \bigl(\E\|\varphi(X_0)\|^2\bigr)^{1/2}<\infty.
\]
Under the source paper's measurability assumptions, $(F_t)$ is therefore a
Bochner-integrable strictly stationary $S_1(\Hh)$-valued sequence.
Pop-Stojanovic's Corollary~2 applies, and ergodicity identifies the limit as
$\E F_0=C(\eta)$. Consequently
\[
 \|C_n(\eta)-C(\eta)\|_{S_1(\Hh)}\longrightarrow 0
 \qquad\text{almost surely}.
\]
This remains true when $\Hh$ is infinite-dimensional, so reflexivity is not
necessary for the convergence in question.  Moreover,
$\|A\|_{S_p}\leq\|A\|_{S_1}$ for every $1\leq p\leq\infty$, so the same
statement yields all Schatten-norm conclusions for each fixed lag.

\section*{Provenance and scope}

The relation is agent-identified: the 1972 supporting paper predates and does
not discuss kernel autocovariance operators.  It nevertheless gives the exact
Banach-valued ergodic theorem needed by Remark~11.  The conclusion uses the
source paper's stationarity, ergodicity, strong measurability, and second-moment
assumptions.  Integer lags form a countable set, so a common probability-one
event for all such lags follows by countable intersection if desired.

\begin{thebibliography}{9}

\bibitem{MKSK}
M. Mollenhauer, S. Klus, C. Sch\"utte, and P. Koltai,
\emph{Kernel Autocovariance Operators of Stationary Processes: Estimation
and Convergence}, arXiv:2004.00891v2 (2022).

\bibitem{Pop}
Z. R. Pop-Stojanovic,
\emph{On ergodic theorem for a Banach valued random sequence},
J. Austral. Math. Soc. \textbf{13} (1972), 501--507,
\href{https://doi.org/10.1017/S1446788700009253}{doi:10.1017/S1446788700009253}.

\end{thebibliography}

\end{document}
